\documentclass[a4paper, 12pt]{article}
\usepackage[english, russian]{babel}
\usepackage[T2A]{fontenc}
\usepackage[utf8]{inputenc}
\usepackage{amsmath}
\usepackage{amssymb}
\usepackage{booktabs}
\usepackage{array}
\usepackage{caption}
\usepackage{lscape} % Для альбомной ориентации
\captionsetup[table]{justification=centering}

\begin{document}
	 % Начало альбомной ориентации
	 
	\begin{table}[h!]
		\centering
		\caption{Сводка решений уравнения Шрёдингера для различных потенциалов}
		\label{tab:potential_solutions}
		\begin{tabular}{>{\raggedright\arraybackslash}p{3.5cm} >{\raggedright\arraybackslash}p{4.5cm} >{\raggedright\arraybackslash}p{4.5cm} >{\raggedright\arraybackslash}p{4cm}}
			\toprule
			\textbf{Потенциал} &
			\textbf{Уравнение Шрёдингера} &
			\textbf{Волновая функция} &
			\textbf{Энергетический спектр} \\
			\midrule
			
			\textbf{Бесконечно глубокая яма ширины $L$}
			&
			$V(x) = \begin{cases}
				0, & 0 < x < L \\
				\infty, & x \leq 0 \text{ или } x \geq L
			\end{cases}$ 
			$-\frac{\hbar^2}{2m}\frac{d^2\psi}{dx^2} = E\psi$
			&
			$\psi_n(x) = \sqrt{\frac{2}{L}}\sin\left(\frac{n\pi x}{L}\right)$ 
			$n = 1, 2, 3, \ldots$ 
			&
			$E_n = \frac{\hbar^2\pi^2 n^2}{2mL^2}$ 
			$n = 1, 2, 3, \ldots$ \\
			\midrule
			
			\textbf{Дельта-потенциал}
			&
			$V(x) = -\alpha\delta(x)$, $\alpha > 0$ \\
			$-\frac{\hbar^2}{2m}\frac{d^2\psi}{dx^2} - \alpha\delta(x)\psi = E\psi$
			&
			$\psi(x) = \sqrt{\kappa}e^{-\kappa|x|}$ \\
			$\kappa = \frac{m\alpha}{\hbar^2}$ \\
			\textit{Одно связанное состояние}
			&
			$E = -\frac{m\alpha^2}{2\hbar^2}$ \\
			\textit{Одно связанное состояние} \\
			\midrule
			
			\textbf{Гармонический осциллятор}
			&
			$V(x) = \frac{kx^2}{2}$ \\
			$-\frac{\hbar^2}{2m}\frac{d^2\psi}{dx^2} + \frac{kx^2}{2}\psi = E\psi$
			&
			$\psi_n(x) = \left(\frac{m\omega}{\pi\hbar}\right)^{1/4}\frac{1}{\sqrt{2^n n!}}H_n(\xi)e^{-\xi^2/2}$ \\
			$\xi = \sqrt{\frac{m\omega}{\hbar}}x$, $\omega = \sqrt{\frac{k}{m}}$ \\
			$H_n$ — полином Эрмита
			&
			$E_n = \hbar\omega\left(n + \frac{1}{2}\right)$ \\
			$n = 0, 1, 2, \ldots$ \\
			\midrule
			
			\textbf{Линейный потенциал}
			&
			$V(x) = F_0x$ \\
			$-\frac{\hbar^2}{2m}\frac{d^2\psi}{dx^2} + F_0x\psi = E\psi$
			&
			$\psi_n(x) = C_n\mathrm{Ai}\left(\frac{x}{x_0} - z_n\right)$ \\
			$x_0 = \left(\frac{\hbar^2}{2mF_0}\right)^{1/3}$ \\
			$C_n$ — нормировочная постоянная
			&
			$E_n = \left(\frac{\hbar^2 F_0^2}{2m}\right)^{1/3}z_n$ \\
			$z_n$ — нули функции Эйри \\
			$z_1 \approx 2.338$, $z_2 \approx 4.088$, \ldots \\
			\bottomrule
		\end{tabular}
	\end{table}% Конец альбомной ориентации

	\noindent \textbf{Примечания:}
	\begin{itemize}
		\item Для бесконечно глубокой ямы: $\psi(0) = \psi(L) = 0$.
		\item Для дельта-потенциала: разрыв производной $\psi'(0^+) - \psi'(0^-) = -\frac{2m\alpha}{\hbar^2}\psi(0)$.
		\item Для гармонического осциллятора: $\omega = \sqrt{\frac{k}{m}}$.
		\item Для линейного потенциала: $\mathrm{Ai}(z)$ — функция Эйри, $z_n$ — её нули.
	\end{itemize}
	
\end{document}
